\documentclass[12pt]{article}
\usepackage{booktabs}
\usepackage{caption}
\usepackage{longtable}
\usepackage{array}
\usepackage[utf8]{inputenc}
\usepackage[margin=0.5in, paperwidth=14in, paperheight=10in]{geometry}
\pagestyle{empty}
\begin{document}
\setcounter{table}{3}
\begin{table}[htbp]
\centering
\large

\caption{Presencia de industrias por departamento (Parte 3 de 3)}
\begin{tabular}[t]{llll}
\toprule
industria & SANTANDER & TOLIMA & VALLE DEL CAUCA\\
\midrule
Alimentos y bebidas & Sí & No & Sí\\
Curtido de cuero y calzado & Sí & No & No\\
Madera y muebles & No & No & Sí\\
Minerales no metálicos & No & No & Sí\\
Papel e imprentas & No & No & Sí\\
\addlinespace
Productos metálicos & No & No & Sí\\
Sustancias y productos químicos, farmacéuticos, de caucho y plástico & Sí & No & Sí\\
Textiles y confecciones & Sí & No & No\\
Textiles, confecciones y cuero & No & No & Sí\\
Total & Sí & Sí & Sí\\
\bottomrule
\end{tabular}
\end{table}
\end{document}
